% ARCHIVED at v3.89 (2026-09-24) from chapters/06_results.tex -- author: Table 6.17 pillars row "over the seeds; FM not flown? that is not the conclusion from the uav-pillars"; 6.4.3 "CI-MeanFM ~ FM ~ MeanFM ~ diffusion? we have proved the pareto optimal" -- the v3.68d ordering sentence contradicted the D3IL-avoiding dominance and is withdrawn.
% Kept verbatim for rollback; not part of the build.

% =============================================================================
% =============================================================================
% 🔒 RECONCILED IN ONE PASS (v3.88, 2026-09-24, author: "6.4 Comparison across Environments at least could
% show all the things besides the scurve now"). The v3.60 lock stands in its rule: this section summarises
% 6.1, 6.2 and 6.3, and every number in it restates a number from those three. Do not "fix" a cell here on
% its own -- if a number disagrees with its source section, the source section is what to check, and the
% section is reconciled again in one pass. UAV-s-curve joins when R44b and R44c land. The v3.60--v3.87 text
% is in withheld/20260924_v3.88_archive/section_6.4_before_v3.88.tex.
% =============================================================================
\section{Comparison across Environments}\label{sec:res:summary}
% =============================================================================

This section recaps the preceding ones and adds no new result. It answers, per environment, the first two
research questions (\autoref{sec:intro:rqs}): how the flow-based models compare with the diffusion baseline
and with one another, and whether endpoint projection beats the per-step projection of \ac{DPCC}.
\hole{UAV-s-curve joins this section when its projected cells and its controller comparison land (R44b,
R44c); until then it is left out of the tables below.}

\subsection{Generative Models}\label{sec:res:summary:models}

Wherever both were run, the flow-based models reach the baseline's success at a small fraction of its time
(\autoref{tab:summary-models}). On the two D3IL tasks their best configurations also match or improve on its
constraint satisfaction; on UAV-corridor they do not, and the baseline keeps the fewest violating steps.
Which of the three objectives is used matters on D3IL-aligning alone, where analytic average-velocity
matching is ahead of the others. On D3IL-avoiding the budget and the selection rule set their order, on
UAV-pillars the two average-velocity models keep that benchmark's order within the spread over the seeds,
and on UAV-corridor, whose demonstrations are straight lines, their flights cannot be told apart.

\begin{table}[htpb]
  \caption[Generative models across environments]{The generative models per environment, as found in
  Sections~\ref{sec:res:state}--\ref{sec:res:uav}: against the diffusion baseline, and between the three
  flow-based objectives. Temporal U-Net of 3.96--4.0\,M parameters throughout; D3IL-aligning on
  training-set contexts. UAV-s-curve is pending.}
  \label{tab:summary-models}
  \centering
  \small
  \begin{tabularx}{\linewidth}{@{}p{2.2cm} L L@{}}
    \toprule
    Environment & Against the baseline & Between the objectives \\
    \midrule
    D3IL-avoiding & CI-MeanFM and FM at $\nfe=1$ dominate it: success with constraint satisfaction $1.000$ held,
                    $59.2$ and $67.0$ control steps against $70.1$, $18.1$ and $17.3$\,ms against $553.4$
                  & the budget and the selection rule set the order; MeanFM at $\nfe=1$ trades one episode in
                    thirty ($0.967$) \\
    D3IL-aligning & MeanFM under per-step projection dominates the baseline's best projected row: nine
                    violation-free contexts against four, $0.179$ against $0.462$\,m, $266$ against $809$\,ms
                  & MeanFM ahead: without projection at $\nfe=20$ it closes $84\,\%$ of the distance, against
                    $14$ (CI-MeanFM), $10$ (FM) and $-2\,\%$ (baseline) \\
    UAV-pillars   & not flown
                  & CI-MeanFM one flight in thirty ahead of MeanFM at both budgets, as on D3IL-avoiding, within
                    the spread over the seeds; FM not flown \\
    UAV-corridor  & every flight reaches the end of the corridor, at a twenty-fourth (tilt) and a ninth (hump)
                    of its time per action; more violating steps, at best $1.4$ and $0.2$ against $0.8$ and $0.0$
                  & none: the same flights without projection, within $1.6$ violating steps per flight at a
                    shared budget with it \\
    \bottomrule
  \end{tabularx}
\end{table}

\subsection{Projection Methods}\label{sec:res:constraints}

Endpoint projection beats the per-step projection of \ac{DPCC} where the task is run at a budget that gives
it steps to guide, and has nothing to win where one or two evaluations solve the task
(\autoref{tab:summary-projection}). On D3IL-aligning, at twenty evaluations, it keeps all ten contexts free
of violations. On UAV-corridor it overtakes the per-step projector over the hump at twenty evaluations with a
single candidate, and under the tilt it leaves fewer violating steps only at more time. On D3IL-avoiding it
is level or ahead at three and ten evaluations, but the task is solved at one, where it has no guiding step,
and UAV-pillars flies that operating point under per-step projection only.

\begin{table}[htpb]
  \caption[Projection methods across environments]{Endpoint against per-step projection per environment, as
  found in Sections~\ref{sec:res:state}--\ref{sec:res:uav}. \emph{Budget}: the network evaluations at which
  the environment is operated or its projectors compared. UAV-s-curve is pending.}
  \label{tab:summary-projection}
  \centering
  \small
  \begin{tabularx}{\linewidth}{@{}p{2.2cm} r L@{}}
    \toprule
    Environment & Budget & Endpoint against per-step projection \\
    \midrule
    D3IL-avoiding & $1$      & no guiding step at the operating budget; at $3$ and $10$, with matched candidates,
                               level or ahead on the constraint and cheaper on every model except FM at $3$ \\
    D3IL-aligning & $20$     & ahead: all ten contexts free of violations and the box moved in eight, at the time
                               per control step of per-step projection; more contexts violation-free in seven of
                               nine comparisons \\
    UAV-pillars   & $1$, $2$ & per-step projection only: no violating step on any flight; the collisions are with
                               pillars outside the constraint sets \\
    UAV-corridor  & $1$--$20$ & ahead only over the hump at $20$ with one candidate, $0.2$ against $1.8$ violating
                               steps at $291$ against $355$\,ms; under the tilt at $20$ fewer violating steps at
                               more time; at $3$ and $5$ as many or more \\
    \bottomrule
  \end{tabularx}
\end{table}

\subsection{Conclusion}\label{sec:res:summary:conclusion}

% v3.68d (author, 2026-09-23): the model orderings beside the table; v3.88: the quadrotor scenes filled in.
Read across the environments, the flow-based objectives separate on the outcome on D3IL-aligning alone. On
D3IL-avoiding they do not separate on the outcome --- CI-MeanFM $\approx$ FM $\approx$ MeanFM $\approx$
diffusion --- and separate on the price by a factor of thirty in favour of the flow-based models. On
D3IL-aligning they separate on the outcome, MeanFM $\gg$ CI-MeanFM $>$ FM $\approx$ diffusion. On UAV-pillars
the two average-velocity models keep their D3IL-avoiding order, CI-MeanFM one flight in thirty ahead. On
UAV-corridor, whose demonstrations are straight lines, the flow-based objectives do not separate, MeanFM
$\approx$ CI-MeanFM $\approx$ FM, and the baseline keeps the fewest violating steps at nine to twenty-four
times their price. The projector orders by the budget the task needs: per-step projection where one or two
evaluations solve it, since endpoint projection has nothing to guide there, and endpoint projection at twenty
evaluations, where it wins --- on D3IL-aligning, and over the hump of UAV-corridor.
\autoref{tab:summary-combinations} collects the combination each environment supports.

\begin{table}[htpb]
  \caption[Selected model--projection combinations]{The supported combination per environment. UAV-s-curve
  is pending.}
  \label{tab:summary-combinations}
  \centering
  \small
  \begin{tabularx}{\linewidth}{@{}p{2.2cm} L L@{}}
    \toprule
    Environment & Combination & Result \\
    \midrule
    D3IL-avoiding
      & MeanFM or CI-MeanFM, $\nfe=1$, per-step projection, temporal consistency
      & S\&C $0.967$ and $1.000$ at $18.3$ and $18.1$\,ms per step; five seeds, the protocol of \ac{DPCC} \\
    D3IL-aligning
      & MeanFM, $\nfe=20$, endpoint projection, random selection
      & $10/10$ contexts free of violations, box moved in $8/10$ \\
    UAV-pillars
      & the D3IL-avoiding configurations: MeanFM or CI-MeanFM, $\nfe=1$ or $2$, per-step projection, temporal
        consistency
      & no violating step; S\&C $0.600$ to $0.700$, every failure a collision with a pillar outside the
        constraint sets \\
    UAV-corridor
      & any flow-based objective, the budget chosen for the trade
      & violating steps per flight from $10.2$ at $27.4$\,ms ($\nfe=1$, per-step) to $1.4$ at $416.9$\,ms
        ($\nfe=20$, endpoint) under the tilt, from $8.0$ at $71.9$\,ms ($\nfe=3$, endpoint) to $0.2$ at
        $291.0$\,ms ($\nfe=20$, endpoint, one candidate) over the hump; every flight reaches the end \\
    \bottomrule
  \end{tabularx}
\end{table}

